Метод предложен бельгийским математиком Виктором Д’Ондтом в 1878 г. и известен также как <<метод делителей $1, 2, 3, \ldots$>>. Он обладает свойством благоприятствовать крупным партиям (завышать их представительство относительно точной пропорции) \cite{balfair,gallprop}.
\subsection{Алгоритм метода Д'Ондта}
\begin{enumerate}
    \item Для каждой партии $i$ вычисляется текущее значение
    \begin{equation*}
        D^{(k)}_i = \frac{V_i}{s_i+1}
    \end{equation*}
    \item Мандат получает та партия, у которой значение 
$ D^{(k)}_i$ максимально. При равенстве решает жребий или дополнительное правило.
\item У победившей партии $s_i$ увеличивается на 1, и алгоритм повторяется, пока все $S$ мандатов не будут распределены.
\end{enumerate}

Крупные партии получают мандаты чаще, чем предписывает строгая пропорция, потому что их голоса делятся на меньшие числа на ранних шагах. Тем временем мелкие партии с малым $V_i$ могут вообще не получить ни одного мандата, даже если их доля голосов превышала бы квоту $Q$, но не достигает выигрышных значений по сравнению с уже поделёнными голосами лидеров \cite{puke}.